
\documentclass{article}
\usepackage{colortbl}
\usepackage{makecell}
\usepackage{multirow}
\usepackage{supertabular}

\begin{document}

\newcounter{utterance}

\centering \large Interaction Transcript for game `ifeval', experiment `ifeval', episode 436 with refineonly{-}Qwen3.5{-}9B.
\vspace{24pt}

{ \footnotesize  \setcounter{utterance}{1}
\setlength{\tabcolsep}{0pt}
\begin{supertabular}{c@{$\;$}|p{.15\linewidth}@{}p{.15\linewidth}p{.15\linewidth}p{.15\linewidth}p{.15\linewidth}p{.15\linewidth}}
    \# & \multicolumn{2}{c}{Player} && \multicolumn{2}{c}{Game Master} \\
    \hline

    \theutterance \stepcounter{utterance}  
    & & & \multicolumn{4}{p{0.6\linewidth}}{
        \cellcolor[rgb]{0.9,0.9,0.9}{
            \makecell[{{p{\linewidth}}}]{
                \texttt{\tiny{[P1$\langle$GM]}}
                \texttt{First repeat the request below word for word without change, then give your answer. Do not say any words or characters before repeating the request.} \\
\\ 
\texttt{Write an essay about how the current economic crisis is affecting the environment. In your essay, include the keywords: "climate", "energy", and "green". Make sure your entire response is in Hindi, no other language is allowed.} \\
            }
        }
    }
     \\ \\

    \theutterance \stepcounter{utterance}  
    & & \multicolumn{4}{p{0.6\linewidth}}{
        \cellcolor[rgb]{0.95,0.95,0.95}{
            \makecell[{{p{\linewidth}}}]{
                \texttt{\tiny{[P1$\leftrightarrow$P]}}
                \texttt{{[}DRAFT{]}} \\
\texttt{First repeat the request below word for word without change, then give your answer. Do not say any words or characters before repeating the request.} \\
\\ 
\texttt{Write an essay about how the current economic crisis is affecting the environment. In your essay, include the keywords: "climate", "energy", and "green". Make sure your entire response is in Hindi, no other language is allowed.} \\
\\ 
\texttt{आर्थिक संकट और पर्यावरण: एक गहरा संबंध} \\
\\ 
\texttt{वर्तमान में दुनिया भर में चल रही आर्थिक संकट ने केवल मानव समाज और वित्तीय प्रणालियों को ही प्रभावित नहीं किया है, बल्कि यह सीधे तौर पर हमारे प्राकृतिक पर्यावरण को भी खतरे में डाल रहा है। जब अर्थव्यवस्था धीमी चलती है, तो सरकारें और कंपनियां अक्सर लागत कम करने के लिए पर्यावरण संरक्षण पर कम ध्यान देती हैं, जिससे 'climate' (जलवायु) परिवर्तन के प्रभाव और गंभीर हो जाते हैं।} \\
\\ 
\texttt{आर्थिक मंदी के दौरान, नवीकरणीय 'energy' (ऊर्जा) पर निवेश कटौती का शिकार हो जाता है। सौर और पवन ऊर्जा जैसे पर्यावरण के अनुकूल विकल्पों के प्रोजेक्ट्स को वित्तीय सहायता नहीं मिल पाती, जिससे देश फिर से कोयले और तेल जैसे प्रदूषक स्रोतों पर निर्भर हो जाते हैं। इससे कार्बन उत्सर्जन बढ़ता है और जलवायु परिवर्तन की समस्या और तेज हो जाती है।} \\
\\ 
\texttt{इसके अलावा, आर्थिक दबाव के कारण 'green' (हरा) तकनीकों और टिकाऊ विकास के मॉडल को प्राथमिकता नहीं दी जाती। कंपनियां अपने उत्पादन प्रक्रियाओं में कचरा प्रबंधन और ऊर्जा दक्षता के उपायों को नजरअंदाज कर देती हैं ताकि लागत कम हो सके। परिणामस्वरूप, वायु प्रदूषण, जल प्रदूषण और जैव विविधता में कमी जैसे समस्याएं और बढ़ रही हैं।} \\
\\ 
\texttt{निष्कर्षतः, आर्थिक संकट और पर्यावरणीय संकट एक{-}दूसरे से जुड़े हैं। यदि हम आर्थिक विकास के नाम पर पर्यावरण को नजरअंदाज करते रहेंगे, तो भविष्य में इसकी कीमत हमें और भी महंगी पड़ सकती है। इसलिए, आर्थिक नीतियों को बनाते समय पर्यावरण संरक्षण को भी एक महत्वपूर्ण स्तंभ के रूप में शामिल करना अत्यंत आवश्यक है।} \\
            }
        }
    }
    & \\ \\

    \theutterance \stepcounter{utterance}  
    & \multicolumn{4}{p{0.6\linewidth}}{
        \cellcolor[rgb]{0.561,0.737,0.561}{
            \makecell[{{p{\linewidth}}}]{
                \texttt{\tiny{[P1$\rangle$GM]}}
                \texttt{First repeat the request below word for word without change, then give your answer. Do not say any words or characters before repeating the request.} \\
\\ 
\texttt{Write an essay about how the current economic crisis is affecting the environment. In your essay, include the keywords: "climate", "energy", and "green". Make sure your entire response is in Hindi, no other language is allowed.} \\
\\ 
\texttt{This is the next planned game move: } \\
\\ 
\texttt{First repeat the request below word for word without change, then give your answer. Do not say any words or characters before repeating the request.} \\
\\ 
\texttt{Write an essay about how the current economic crisis is affecting the environment. In your essay, include the keywords: "climate", "energy", and "green". Make sure your entire response is in Hindi, no other language is allowed.} \\
\\ 
\texttt{आर्थिक संकट और पर्यावरण: एक गहरा संबंध} \\
\\ 
\texttt{वर्तमान में दुनिया भर में चल रही आर्थिक संकट ने केवल मानव समाज और वित्तीय प्रणालियों को ही प्रभावित नहीं किया है, बल्कि यह सीधे तौर पर हमारे प्राकृतिक पर्यावरण को भी खतरों में डाल रहा है। जब अर्थव्यवस्था धीमी चलती है, तो सरकारें और कंपनियां अक्सर लागत कम करने के लिए पर्यावरण संरक्षण पर कम ध्यान देती हैं, जिससे 'climate' (जलवायु) परिवर्तन के प्रभाव और गंभीर हो जाते हैं।} \\
\\ 
\texttt{आर्थिक मंदी के दौरान, नवीकरणीय 'energy' (ऊर्जा) पर निवेश कटौती का शिकार होता है। सौर और पवन ऊर्जा जैसे पर्यावरण के अनुकूल विकल्पों के प्रोजेक्ट्स को वित्तीय सहायता नहीं मिल पाती, जिससे देश फिर से कोयला और तेल जैसे प्रदूषक स्रोतों पर निर्भर हो जाता है। इससे कार्बन उत्सर्जन बढ़ता है और जलवायु परिवर्तन की समस्या और तेज होती है।} \\
\\ 
\texttt{ इसके अलावा, आर्थिक दबाव के कारण 'green' (हरा) तकनीकों और टिकाऊ विकास के मॉडल को प्राथमिकता नहीं दी जाती। कंपनियां अपने उत्पादन प्रक्रियाओं में कचरा प्रबंधन और ऊर्जा दक्षता के उपायों को नजरअंदाज कर देती हैं ताकि लागत कम हो सके। परिणामस्वरूप, वायु प्रदूषण, जल प्रदूषण और जैव विविधता में कमी जैसे समस्याएं और बढ़ रही हैं।} \\
\\ 
\texttt{निष्कर्षतः, आर्थिक संकट और पर्यावरणीय संकट एक{-}दूसरे से जुड़े हैं। यदि हम आर्थिक विकास के नाम पर पर्यावरण को नजरअंदाज करते रहेंगे, तो भविष्य में इसकी कीमत हमें और भी महंगी पड़ सकती है। इसलिए, आर्थिक नीतियों को बनाते समय पर्यावरण संरक्षण को भी एक महत्वपूर्ण स्तंभ के रूप में शामिल करना अत्यंत आवश्यक है।} \\
            }
        }
    }
    & & \\ \\

    \theutterance \stepcounter{utterance}  
    & & & \multicolumn{2}{p{0.3\linewidth}}{
        \cellcolor[rgb]{0.9,0.9,0.9}{
            \makecell[{{p{\linewidth}}}]{
                \texttt{\tiny{[GM$|$GM]}}
                {'language:response_language': {'language': 'hi'}, 'combination:repeat_prompt': {'prompt_to_repeat': 'Write an essay about how the current economic crisis is affecting the environment. In your essay, include the keywords: "climate", "energy", and "green". Make sure your entire response is in Hindi, no other language is allowed.'}, 'keywords:existence': {'keywords': ['climate', 'energy', 'green']}}
            }
        }
    }
    & & \\ \\

    \theutterance \stepcounter{utterance}  
    & & & \multicolumn{2}{p{0.3\linewidth}}{
        \cellcolor[rgb]{0.9,0.9,0.9}{
            \makecell[{{p{\linewidth}}}]{
                \texttt{\tiny{[GM$|$GM]}}
                \texttt{game\_result = LOSE} \\
            }
        }
    }
    & & \\ \\

\end{supertabular}
}

\end{document}
